% Generated draft: requires mathematical and source review.
\documentclass[11pt]{article}
\usepackage[margin=25mm]{geometry}
\usepackage{fontspec}
\setmainfont{Latin Modern Roman}
\usepackage{amsmath,amssymb,mathtools,unicode-math}
\setmathfont{Latin Modern Math}
\usepackage{enumitem,needspace,xurl,hyperref}
\setlength{\parindent}{0pt}
\setlength{\parskip}{5pt}
\setlength{\emergencystretch}{3em}
\newcommand{\N}{\mathbb{N}}
\newcommand{\Z}{\mathbb{Z}}
\newcommand{\Q}{\mathbb{Q}}
\newcommand{\R}{\mathbb{R}}
\newcommand{\C}{\mathbb{C}}
\newcommand{\F}{\mathbb{F}}
\newcommand{\A}{\mathbb{A}}
\newcommand{\PP}{\mathbb{P}}
\newcommand{\E}{\mathbb{E}}
\newcommand{\eps}{\varepsilon}
\newcommand{\dd}{\,\mathrm{d}}
\begin{document}
\section*{UnsolvedMath ID: 30001691}
\textbf{Source problem code:} catalogue entry
\subsection{Definitions and mathematical statement}
The supplied question is: ``Can unknown non-separated inclusions be identified from the available inverse-problem data?'' Here an inclusion is an unknown region or configuration in the underlying inverse problem. The supplied record does not specify the ambient domain, governing equation, admissible configurations, physical parameters, or measured data. In particular, it does not define ``non-separated'': no condition concerning disjointness, contact, overlap, or distance between inclusions is supplied.

A precise mathematical formulation therefore remains unresolved. To describe the missing structure, let \(\mathcal A\) denote a specified class of admissible inclusion configurations, let \(\mathcal Y\) denote a specified space of observations, and let \(F\colon\mathcal A\to\mathcal Y\) denote the corresponding forward observation map. These are placeholders, not objects determined by the supplied record. If ``identified'' means unique determination from exact data, and if \(\sim\) is the specified equivalence relation describing when two configurations count as the same, the corresponding question would be whether
\[
\forall D_1,D_2\in\mathcal A,\qquad F(D_1)=F(D_2)\ \Longrightarrow\ D_1\sim D_2.
\]
This is only a conditional formulation of one possible interpretation. The record does not establish whether the intended target is uniqueness, a reconstruction procedure, stability under measurement error, or some combination of these. It also does not specify whether the number of inclusions is known, which coefficients are known, what boundary conditions apply, or how much measurement data is available. Consequently, neither this conditional formulation nor any particular inverse boundary-value problem can be asserted as the exact intended mathematical statement without further source context.
\subsection{Short English statement}
The question asks whether inverse-problem measurements determine unknown non-separated inclusions. The supplied record leaves the forward model, measurements, admissible inclusions, meaning of separation, and meaning of identification unspecified, so a precise problem requires further source context.
\subsection{Sources}
\begin{itemize}
\item Local registry: \texttt{data/unsolvedmath/problems.json}, permanent UnsolvedMath ID \(30001691\), source code \texttt{catalogue entry}; supplied title: \emph{Inverse Identification of Nonseparated Inclusions}. The record belongs to the Oberwolfach Reports Open Problems collection and records proposed year 2011; no individual proposer is supplied.
\item Supplied report citation: \emph{Mini-Workshop: Nonlinear Least Squares in Shape Identification Problems} (2011), Oberwolfach Reports, pp.~201--239, DOI 10.4171/owr/2011/05. Source URL and literature-source label ``original OWR report'': \url{https://doi.org/10.4171/owr/2011/05}. The record describes these reports as documenting mathematical workshops and their open-problem discussions. No location within the report more specific than this page range is supplied.
\item Additional supplied literature source, labeled ``planar conductivity-inclusion recovery'': \url{https://doi.org/10.1137/22M1522395}. No authors, title, or detailed theorem statement for this reference are supplied.
\item Supplied literature assessment, checked on exactly \texttt{2026-08-21}: the status is labeled open; the assessment says that the report explicitly labels recovery of unknown non-separated inclusions open and that later work located by the curation treats restricted inclusion classes rather than the general formulation. This is supplied dated curation derived from the Oberwolfach Reports material, not an independently verified assessment. Time-sensitive status and cited claims require independent verification.
\item The registry marks the statement status as ``exact'' and reports that its mathematical content was checked against the official report, with notation, context, and question form normalized without changing the source-level claim. These are supplied verification claims; no source was visited or independently checked in preparing this draft.
\end{itemize}
\end{document}
